\section{Дерево вывода}
\label{sect:DerivTree}

В некоторых случаях не достаточно знать порядок применения правил.
Необходимо структурное представление вывода цепочки в грамматике.
Таким представлением является \emph{дерево вывода}.

\begin{definition}[Дерево вывода]
    Деревом вывода цепочки $\omega$ в грамматике $G = \langle \Sigma, N, S, P \rangle$ называется корневое упорядоченное помеченное дерево, удовлетворяющее следующим свойствам.
    \begin{enumerate}
        \item Метка каждого внутреннего узла~--- нетерминал. При этом метка корня~--- стартовый нетерминал
        \item Метка каждого листа~--- терминал или $\varepsilon$.
        \item В дереве может существовать узел с меткой $N_i$ и сыновьями $M_j \dots M_k$ только тогда, когда в грамматике есть правило вида $N_i \to M_j \dots M_k$.
        \item Крона образует исходную цепочку $\omega$.
    \end{enumerate}
\end{definition}

\begin{example}
    Пусть дана грамматика
    \begin{equation}
        G = \langle \{a,b\}, \{S\}, S, \{S \to a \ S \ b \ S, S \to \varepsilon\} \rangle.\label{eq:grammar}
    \end{equation}
    Дерево вывода цепочки $abab$ в этой грамматике представлено на рисунке~\ref{fig:derivation_tree_example}.

    \begin{marginfigure}
        \begin{center}
            \resizebox{\marginparwidth}{!}{\input{part_02_Foundations/chapter_06_ContextFreeLanguages/figures/tree0.tex}}
        \end{center}
        \caption{Дерева вывода цепочки $abab$ в грамматике~\ref{eq:grammar}}
        \label{fig:derivation_tree_example}
    \end{marginfigure}

\end{example}

\begin{theorem}
    Пусть $G = \langle \Sigma, N, P, S \rangle$~--- КС-грамматика.
    Вывод $S \derives \alpha$, где $\alpha \in (N \cup \Sigma)^*, \alpha \neq \varepsilon$ существует тогда и только тогда, когда существует дерево вывода в грамматике $G$ с кроной $\alpha$.
\end{theorem}
